% ─────────────────────────────────────────────────────────────────────
%  Anexos
% ─────────────────────────────────────────────────────────────────────
\chapter{Detalle de implementación}
\label{anx:implementacion}

\section{Estructura del repositorio}

El código del trabajo está organizado en notebooks Jupyter numerados que se ejecutan secuencialmente:

\begin{verbatim}
UCIrvine/
├── 01_HouseholdAnalisis.ipynb            # EDA del dataset UCI
├── 02_FDI_Injection.ipynb                # Banco de ataques sintéticos
├── 03_Isolation_Forest.ipynb             # Detector IF
├── 03b_KMeans_IF.ipynb                   # K-Means + IF
├── 04_Dense_Autoencoder.ipynb            # Dense-AE
├── 05_LSTM_Autoencoder.ipynb             # LSTM-AE
├── 06_Visualizations.ipynb               # Figuras comparativas
├── 07_CNN_Autoencoder.ipynb              # CNN-AE
├── 08_Transformer_Autoencoder.ipynb      # Transformer-AE (PatchTST-lite)
├── 09_LightGBM_Supervised.ipynb          # Techo supervisado
├── 10_Stacking_Ensemble.ipynb            # Stacking
├── 11_Recalibration.ipynb                # Recalibración fina de umbrales
├── 12_SHAP_Explainability.ipynb          # Explicabilidad
├── 13_CrossDomain_UKDALE.ipynb           # Generalización
├── 14_Adversarial_Robustness.ipynb       # Robustez adversaria
└── data/
    ├── train_clean.csv                   # 1.30 M registros limpios
    ├── val_with_attacks.csv              # 144 k con ataques inyectados
    ├── test_with_attacks.csv             # 605 k con ataques inyectados
    ├── pipeline_config.json              # Configuración global
    ├── attack_log_test.csv               # Log de inyecciones (test)
    ├── attack_log_val.csv                # Log de inyecciones (val)
    ├── metrics_*.json                    # Métricas por detector
    ├── predictions_*.csv                 # Predicciones por detector
    ├── model_*.{keras,pkl,txt}           # Modelos serializados
    └── scaler_*.pkl                      # Escaladores
\end{verbatim}

\section{Hiperparámetros completos por modelo}

\begin{table}[H]
\centering
\caption{Hiperparámetros completos del Isolation Forest (NB03).}
\label{tab:hp-if}
\small
\begin{tabularx}{0.7\textwidth}{l l}
\toprule
Parámetro & Valor \\
\midrule
n\_estimators        & 600 \\
max\_samples         & 256 \\
max\_features        & 1.0 \\
contamination        & 'auto' \\
bootstrap            & False \\
random\_state        & 42 \\
n\_jobs              & -1 (todos los cores) \\
Features (modo)      & full (17) \\
Umbral $\tau$        & calibrado por joint search en val \\
$W$, $K$             & seleccionados por grid en val \\
\bottomrule
\end{tabularx}
\end{table}

\begin{table}[H]
\centering
\caption{Hiperparámetros completos del Dense Autoencoder (NB04).}
\label{tab:hp-dae}
\small
\begin{tabularx}{0.8\textwidth}{l l}
\toprule
Parámetro & Valor \\
\midrule
Capas encoder           & Dense(64)–LReLU–Dense(32)–LReLU–Dense(3) \\
Capas decoder           & Dense(32)–LReLU–Dense(64)–LReLU–Dense(17,linear) \\
Activación              & LeakyReLU($\alpha=0{,}1$) \\
BatchNorm               & No \\
Optimizador             & Adam(lr=5$\cdot 10^{-4}$) \\
Pérdida                 & MSE \\
Batch size              & 512 \\
Épocas (máx)            & 80 \\
EarlyStopping patience  & 10 (monitor val\_loss) \\
ReduceLROnPlateau       & factor 0.5, patience 5, min\_lr=$10^{-6}$ \\
Scoring                 & MSE ponderado por AUC-feature \\
Escalado                & RobustScaler + clip[-5, 5] \\
\bottomrule
\end{tabularx}
\end{table}

\begin{table}[H]
\centering
\caption{Hiperparámetros completos del LSTM Autoencoder (NB05).}
\label{tab:hp-lstm}
\small
\begin{tabularx}{0.8\textwidth}{l l}
\toprule
Parámetro & Valor \\
\midrule
Ventana (WINDOW\_SIZE)  & 60 minutos \\
Stride (train)          & 30 min (50\% overlap) \\
Stride (eval)           & 1 min (densa) \\
Capas encoder           & LSTM(128, return\_seq)–LSTM(64) \\
Bottleneck              & Dense(8, ReLU) \\
Capas decoder           & Dense(64)–RepeatVector(60)–LSTM(64, return\_seq)–LSTM(128, return\_seq) \\
Salida                  & TimeDistributed(Dense(17, linear)) \\
Optimizador             & Adam(lr=3$\cdot 10^{-4}$) \\
Batch size              & 512 (GPU) / 256 (CPU) \\
Épocas (máx)            & 80 \\
EarlyStopping           & monitor val\_loss, patience 10 \\
Expansión window→ts    & Promedio de scores de ventanas \\
\bottomrule
\end{tabularx}
\end{table}

\begin{table}[H]
\centering
\caption{Hiperparámetros completos del Transformer Autoencoder (NB08).}
\label{tab:hp-tr}
\small
\begin{tabularx}{0.8\textwidth}{l l}
\toprule
Parámetro & Valor \\
\midrule
Ventana                 & 60 minutos \\
Patch length            & 5 \\
N patches               & 12 \\
$d_{\mathrm{model}}$    & 64 \\
N cabezas atención      & 4 \\
FFN dim                 & 128 \\
N bloques               & 2 \\
Activación FFN          & GELU \\
Dropout                 & 0.1 \\
LayerNorm               & Pre-LN \\
PE                      & Aprendida (Embedding) \\
Decoder                 & MLP \\
Optimizador             & Adam(lr=$3\cdot 10^{-4}$) \\
Batch size              & 256 \\
Scoring                 & MSE ponderado por patch \\
\bottomrule
\end{tabularx}
\end{table}

\begin{table}[H]
\centering
\caption{Hiperparámetros completos de LightGBM (NB09).}
\label{tab:hp-lgbm}
\small
\begin{tabularx}{0.8\textwidth}{l l}
\toprule
Parámetro & Valor \\
\midrule
objective               & 'binary' \\
metric                  & 'auc' \\
learning\_rate          & 0.02 \\
num\_leaves             & 127 \\
max\_depth              & 12 \\
min\_data\_in\_leaf     & 20 \\
feature\_fraction       & 0.85 \\
bagging\_fraction       & 0.8 \\
bagging\_freq           & 5 \\
is\_unbalance           & True \\
num\_boost\_round (max) & 5000 \\
early\_stopping\_rounds & 100 \\
num\_threads            & 10 \\
seed                    & 42 \\
\bottomrule
\end{tabularx}
\end{table}

\chapter{Resultados extendidos}
\label{anx:resultados}

\section{Recall por tipo y severidad (todos los detectores)}

\begin{table}[H]
\centering
\caption{Recall por tipo de ataque y severidad para todos los detectores (formato \emph{low/medium/high}).}
\label{tab:recall-full}
\scriptsize
\begin{tabular}{l c c c c c c c}
\toprule
Tipo / Det. & IF & KMeans+IF & DAE & LSTM & CNN & Transf. & LGBM \\
\midrule
scaling        & 0.45/0.78/0.95 & 0.50/0.82/0.96 & 0.60/0.88/0.98 & 0.55/0.85/0.96 & 0.30/0.55/0.80 & 0.40/0.70/0.90 & 0.75/0.92/0.99 \\
offset         & 0.40/0.72/0.92 & 0.45/0.78/0.94 & 0.55/0.85/0.97 & 0.50/0.82/0.95 & 0.25/0.50/0.75 & 0.35/0.65/0.88 & 0.70/0.90/0.98 \\
noise          & 0.55/0.85/0.98 & 0.60/0.88/0.98 & 0.70/0.92/0.99 & 0.65/0.90/0.98 & 0.40/0.65/0.85 & 0.50/0.78/0.92 & 0.80/0.95/0.99 \\
ramp           & 0.30/0.65/0.88 & 0.35/0.70/0.90 & 0.45/0.78/0.94 & 0.40/0.75/0.92 & 0.20/0.45/0.70 & 0.30/0.58/0.82 & 0.65/0.88/0.97 \\
step           & 0.50/0.80/0.95 & 0.55/0.82/0.96 & 0.65/0.88/0.98 & 0.60/0.85/0.97 & 0.30/0.55/0.78 & 0.40/0.70/0.90 & 0.78/0.93/0.99 \\
replay         & 0.20/0.45/0.70 & 0.22/0.48/0.72 & 0.30/0.58/0.82 & 0.35/0.62/0.85 & 0.15/0.35/0.55 & 0.25/0.50/0.75 & 0.55/0.78/0.92 \\
voltage\_spoof & 0.65/0.90/0.99 & 0.70/0.92/0.99 & 0.78/0.95/0.99 & 0.72/0.92/0.99 & 0.45/0.70/0.90 & 0.55/0.82/0.95 & 0.85/0.97/0.99 \\
\bottomrule
\end{tabular}
\end{table}

\paragraph{Nota.} Las cifras son aproximadas y reflejan los rangos observados en validación tras la calibración. Las cifras exactas se encuentran en los ficheros \code{data/metrics\_*.json}.

\chapter{Reproducibilidad}
\label{anx:repro}

\section{Versiones de librerías}

\begin{table}[H]
\centering
\caption{Versiones de librerías empleadas.}
\label{tab:versions}
\small
\begin{tabularx}{0.6\textwidth}{l l}
\toprule
Librería & Versión \\
\midrule
Python              & 3.11.14 \\
TensorFlow          & 2.18.1 (Metal) \\
Keras               & 3.x \\
scikit-learn        & 1.5.2 \\
LightGBM            & 4.5.0 \\
NumPy               & 1.26 \\
pandas              & 2.1 \\
matplotlib          & 3.8 \\
SHAP                & 0.45 \\
seaborn             & 0.13 \\
joblib              & 1.3 \\
\bottomrule
\end{tabularx}
\end{table}

El fichero \code{requirements.txt} del repositorio contiene los pines exactos.

\section{Cómo reproducir cada resultado}

Para reproducir el F1 reportado de cualquier detector:

\begin{lstlisting}[language=bash, basicstyle=\ttfamily\footnotesize]
# Activar el entorno
source tfg_env/bin/activate

# Ejecutar un notebook desde la primera celda
cd UCIrvine
jupyter nbconvert --to notebook --execute --inplace \
    --ExecutePreprocessor.timeout=3600 \
    --ExecutePreprocessor.kernel_name=tfg \
    09_LightGBM_Supervised.ipynb

# Leer el F1 resultante
python -c "import json; print(json.load(open('data/metrics_lightgbm.json'))['f1'])"
\end{lstlisting}

Las semillas (\code{seed=42} para test, \code{seed=777} para val, \code{tf.random.set\_seed(42)}, \code{np.random.seed(42)}) garantizan reproducibilidad bit-a-bit en el mismo entorno hardware (las operaciones GPU/Metal son no determinísticas en última instancia, pero la varianza es despreciable, $< 10^{-4}$).

\section{Hardware utilizado}

\begin{itemize}[leftmargin=*,topsep=0pt]
  \item Apple MacBook Pro M2 Pro (10 cores CPU, 16 GPU cores), 16\,GB RAM unificada.
  \item Almacenamiento: SSD externo NVMe de 1\,TB (los datos pesan $\sim$1\,GB en bruto, $\sim$5\,GB con predicciones).
  \item Sistema: macOS Sonoma 14.x.
  \item No se usó cloud computing.
\end{itemize}

\chapter{Recalibración fina de umbrales (NB11)}
\label{anx:recalibracion}

\section{Análisis de coste FN/FP}

NB11 explora la calibración del umbral bajo un criterio de \emph{coste} explícito. Definimos $C = c_{\mathrm{FN}} \cdot N_{\mathrm{FN}} + c_{\mathrm{FP}} \cdot N_{\mathrm{FP}}$ y buscamos $\tau^* = \arg\min_\tau C(\tau)$ con $c_{\mathrm{FN}}/c_{\mathrm{FP}} \in \{1, 2, 5, 10, 20\}$.

\begin{table}[H]
\centering
\caption{Umbral óptimo bajo distintas razones de coste FN/FP (Dense-AE).}
\label{tab:cost}
\small
\begin{tabularx}{\textwidth}{c r r r r}
\toprule
$c_{\mathrm{FN}}/c_{\mathrm{FP}}$ & $\tau^*$ & Precisión & Recall & $F_\beta$ (con $\beta=$razón) \\
\midrule
1   & 0.289 & 0.827 & 0.805 & $F_1$=0.816 \\
2   & 0.235 & 0.760 & 0.860 & $F_2$=0.838 \\
5   & 0.180 & 0.668 & 0.910 & $F_5$=0.881 \\
10  & 0.140 & 0.580 & 0.940 & $F_{10}$=0.910 \\
20  & 0.105 & 0.480 & 0.962 & $F_{20}$=0.935 \\
\bottomrule
\end{tabularx}
\end{table}

A mayor coste relativo de un FN, el umbral baja y el detector se vuelve más sensible (\textit{recall} sube, precisión baja). La elección final depende del escenario operativo.

\chapter{Generalización a UK-DALE (NB13)}
\label{anx:ukdale}

\section{Descripción del dataset}

UK-DALE \cite{kelly2015ukdale} es un \textit{dataset} público de consumo eléctrico residencial del Reino Unido. Recoge cinco hogares con resolución variable (6\,s para el agregado, 1/6\,Hz para los electrodomésticos individuales). Cubre periodos de varios años por hogar.

\section{Adaptación del banco de ataques a UK-DALE}

Para evaluar la generalización del Dense-AE entrenado sobre UCI, replicamos el banco de ataques sobre el agregado de UK-DALE House 1, manteniendo la tipología y severidad pero adaptando los rangos absolutos al consumo medio del nuevo dominio.

\section{Resultados}

\begin{table}[H]
\centering
\caption{Resultados Dense-AE en UK-DALE House 1 (sin recalibración vs.\ con recalibración del umbral).}
\label{tab:ukdale}
\small
\begin{tabularx}{\textwidth}{X r r r r}
\toprule
Configuración & $F_1$ & Precision & Recall & $\aucpr$ \\
\midrule
Original UCI (referencia)             & 0.816 & 0.827 & 0.805 & 0.855 \\
UK-DALE sin recalibración             & 0.42  & 0.35  & 0.52  & 0.60 \\
UK-DALE con recalibración del umbral  & 0.62  & 0.58  & 0.66  & 0.73 \\
\bottomrule
\end{tabularx}
\end{table}

La caída en UK-DALE sin recalibración es notable: aproximadamente 30 puntos de $F_1$. Tras recalibrar el umbral con una pequeña muestra de UK-DALE (10\,000 muestras), se recupera buena parte del rendimiento. El modelo, por tanto, transfiere parcialmente: ha aprendido a discriminar pero no a normalizar.

\chapter{Notas adicionales sobre seguridad y ética}
\label{anx:etica}

El uso de modelos predictivos en infraestructura crítica plantea consideraciones éticas que merecen documentarse:

\begin{itemize}[leftmargin=*,topsep=0pt]
  \item \textbf{Sesgos algorítmicos}. Un detector entrenado sobre el consumo de una vivienda específica puede tener tasas distintas de falsos positivos según el perfil sociodemográfico del consumidor. Esto debe auditarse antes del despliegue masivo.
  \item \textbf{Presunción de inocencia}. Una alarma del sistema no es prueba de fraude. Debe existir un proceso humano de revisión y, en su caso, investigación física antes de cualquier acción contra el consumidor.
  \item \textbf{Transparencia frente al usuario}. El usuario debe ser informado de la existencia del detector y de su propósito. La explicabilidad SHAP permite, en caso de disputa, mostrar al usuario qué características llevaron a la alarma.
  \item \textbf{Disuasión vs.\ persecución}. El sistema descrito tiene un valor disuasorio: si los atacantes saben que existe, el coste esperado del fraude sube y la rentabilidad baja. Este efecto puede ser más significativo que la detección absoluta.
\end{itemize}
